\begin{enumerate}
  \item Écrivons l'équation de la réaction et dressons son tableau d'évolution en effectuant un bilan en quantités de matière, puis en concentrations.
\end{enumerate}

$$
\mathrm{H}-\mathrm{COOH}(a q)+\mathrm{H}_{2} \mathrm{O}=\mathrm{H}-\mathrm{COO}^{-}(a q)+\mathrm{H}_{3} \mathrm{O}^{+}
$$

\begin{center}
\begin{tabular}{|c|c|c|c|c|c|}
\hline
\multirow{2}{*}{$\mathbf{B M}$} & \begin{tabular}{c}
État \\
initial \\
\end{tabular} & $n$ & $\sim$ & 0 & 0 \\
\cline { 2 - 6 }
 & \begin{tabular}{c}
État \\
final \\
\end{tabular} & $n-x_{f}$ & $\sim$ & $x_{f}$ & $x_{f}$ \\
\hline
\multirow{2}{*}{$\mathbf{B C}$} & \begin{tabular}{c}
État \\
initial \\
\end{tabular} & $c$ & $\sim$ & 0 & 0 \\
\cline { 2 - 6 }
 & \begin{tabular}{c}
État \\
final \\
\end{tabular} & $c-\left[\mathrm{H}_{3} \mathrm{O}^{+}\right]_{f}$ & - & $[\mathrm{H}-\mathrm{COO}(a q)]_{f}=\left[\mathrm{H}_{3} \mathrm{O}^{+}\right]_{f}$ & $\left[\mathrm{H}_{3} \mathrm{O}^{+}\right]_{f}$ \\
\hline
\end{tabular}
\end{center}

Le caractère conducteur de la solution est dû aux ions $\mathrm{H}-\mathrm{COO}^{-}(a q)$ et $\mathrm{H}_{3} \mathrm{O}^{+}$produits par la réaction. La conductivité de la solution est:

$$
\begin{gathered}
\sigma=G \times \frac{\ell}{S}=3,30 \times 10^{-4} \times \frac{10^{-2}}{10^{-4}}=3,30 \times 10^{-2} \mathrm{~S} \cdot \mathrm{~m}^{-1} \\
\sigma=\lambda_{\mathrm{H}_{3} \mathrm{O}^{+}} \times\left[\mathrm{H}_{3} \mathrm{O}^{+}\right]_{f}+\lambda_{\mathrm{H}-\mathrm{COO}} \times[\mathrm{H}-\mathrm{COO}-(a q)]_{f}=\left(\lambda_{\mathrm{H}_{3} \mathrm{O}^{+}}+\lambda_{\mathrm{H}-\mathrm{COO}}\right) \times\left[\mathrm{H}_{3} \mathrm{O}^{+}\right]_{f}
\end{gathered}
$$

D'où : $\left[\mathbf{H}_{\mathbf{3}} \mathbf{O}^{+}\right]_{f}=\frac{\sigma}{\lambda_{\mathrm{H}_{3} \mathrm{O}^{+}}+\lambda_{\mathrm{H}-\mathrm{COO}}}=\frac{3,30 \times 10^{-2}}{(3,50+0,55) \times 10^{-2}}=\mathbf{0 , 8 1 5 ~ m o l} \cdot \mathbf{m}^{-\mathbf{3}}$ ou $\left[\mathrm{H}_{3} \mathrm{O}^{+}\right]_{f}=8,15 \times 10^{-4} \mathrm{~mol} \cdot \mathrm{~L}^{-1}$.\\
Le taux d'avancement final de la réaction est : $\tau=\frac{x_{f}}{n}=\frac{\left[\mathrm{H}_{3} \mathrm{O}^{+}\right]_{f}}{c}=\frac{8,15 \times 10^{-4}}{5,0 \times 10^{-3}}$.

$$
\tau=0,163 \text {, soit } 16,3 \%
$$

\begin{enumerate}
  \setcounter{enumi}{1}
  \item La constante d'équilibre $K$ associée à l'équation de la réaction s'écrit:
\end{enumerate}

$$
\begin{array}{r}
K=Q_{r e q}=\frac{[\mathrm{H}-\mathrm{COO}]_{f} \times\left[\mathrm{H}_{3} \mathrm{O}^{+}\right]_{f}}{[\mathrm{H}-\mathrm{COOH}]_{f}} \\
K=\frac{\left[\mathrm{H}_{3} \mathrm{O}^{+}\right]^{2}}{c-\left[\mathrm{H}_{3} \mathrm{O}^{+}\right]}=\frac{\left(8,15 \times 10^{-4}\right)^{2}}{5,0 \times 10^{-3}-8,15 \times 10^{-4}} ; \boldsymbol{K = 1 , 5 9 \times 1 0 ^ { - 4 }}
\end{array}
$$

